\documentclass[11pt,a4paper]{article}
\usepackage[margin=1in]{geometry}
\usepackage{amsmath,amssymb,amsthm}
\usepackage[hidelinks]{hyperref}

\newtheorem{theorem}{Theorem}
\newtheorem{proposition}[theorem]{Proposition}
\newtheorem{conjecture}[theorem]{Conjecture}

\title{Disproof of Li's Conjecture on Twin Practical Numbers of the Form $q^C$ and $q^C + 2$}
\author{Salem Levinstein}
\date{November 27, 2026}

\begin{document}
\maketitle

\begin{abstract}
Runbo Li (\emph{Notes on Number Theory and Discrete Mathematics} 32(3), 2026, pp.\ 544--555) conjectured that for every positive integer $C$, there exist infinitely many practical numbers $q$ such that both $q^C$ and $q^C + 2$ are practical numbers (Conjecture~1.1). We disprove this conjecture for $C = 30$ and, more generally, for every positive multiple $C$ of $30$: for every positive integer $q$, the number $q^{30} + 2$ is never practical. The proof is an elementary local obstruction modulo $4, 9, 5, 7, 11, 13$ showing that $13$ (or $2$ when $q$ is odd) cannot be written as a sum of distinct divisors of $q^{30} + 2$. We further prove that this local obstruction holds for an exponent $C \ge 1$ if and only if $30 \mid C$, establishing that $C = 30$ is the exact minimal modulus of the obstruction, and contrast this with explicit practical triples $(q, q^C, q^C + 2)$ for all $1 \le C \le 6$. All results are formally machine-checked in Lean~4 (\texttt{v4.33.1} + Mathlib, standard axioms only).
\end{abstract}

\section{Introduction}

A positive integer $n \ge 1$ is called \emph{practical} if every integer $m$ with $1 \le m \le n$ can be written as a sum of distinct positive divisors of $n$. Practical numbers were introduced by Srinivasan \cite{Srinivasan1948} and completely characterized via their prime factorization by Stewart \cite{Stewart1954} and Sierpi\'{n}ski \cite{Sierpinski1955}: $n = 1$ is practical, and an integer $n = p_1^{a_1} \cdots p_k^{a_k} > 1$ with primes $p_1 < \cdots < p_k$ is practical if and only if $p_1 = 2$ and $p_i \le \sigma(p_1^{a_1} \cdots p_{i-1}^{a_{i-1}}) + 1$ for every $2 \le i \le k$.

In a recent paper, Runbo Li \cite{Li2026} established the existence of infinitely many practical numbers $q$ with $q^p$ and $q^p + 2$ both practical for every prime $7 \le p < 100$ (extending the cases $C = 1$ due to Margenstern and Melfi and $C = 2, 4$ due to Wang and Sun \cite{WangSun2018}), and proposed the following general conjecture.

\begin{conjecture}[Li \cite{Li2026}, Conjecture~1.1]\label{conj:li}
For any positive integer $C$, there are infinitely many practical numbers $q$ such that $q^C$ and $q^C + 2$ are also practical numbers.
\end{conjecture}

We show that Conjecture~\ref{conj:li} fails in the strongest possible sense whenever $30 \mid C$: not a single positive integer $q$ satisfies the conclusion.

\section{Refutation for All Multiples of $30$}

\begin{theorem}\label{thm:main}
For every positive integer $q$, the integer $q^{30} + 2$ is not practical. Consequently, for every positive integer $C$ divisible by $30$, there is no positive integer $q$ such that $q^C + 2$ is practical, and Conjecture~\ref{conj:li} is false.
\end{theorem}

\begin{proof}
Let $q \ge 1$ be an integer and set $n = q^{30} + 2$.

\medskip\noindent\emph{Case 1: $q$ is odd.}
Then $n \ge 3$ is odd, so every positive divisor of $n$ is odd. In particular, the only positive divisor of $n$ at most $2$ is $1$. Any subset $S$ of positive divisors of $n$ whose sum is $2$ can contain only divisors $\le 2$, so $S \subseteq \{1\}$ and $\sum_{d \in S} d \le 1 < 2$, a contradiction. Thus $n$ is not practical.

\medskip\noindent\emph{Case 2: $q$ is even.}
Then $q \ge 2$, so $n = q^{30} + 2 \ge 2^{30} + 2 > 13$. We examine $n$ modulo six small prime powers:
\begin{itemize}
\item Modulo $4$: since $2 \mid q$, we have $4 \mid q^{30}$, so $n \equiv 2 \pmod{4}$ and hence $4 \nmid n$.
\item Modulo $9$: if $3 \mid q$, then $q^{30} \equiv 0 \pmod{9}$; otherwise $\gcd(q, 9) = 1$, and Euler's theorem ($\varphi(9) = 6$) gives $q^{30} = (q^6)^5 \equiv 1 \pmod{9}$. Thus $n \in \{2, 3\} \pmod{9}$, so $9 \nmid n$.
\item Modulo $5$: if $5 \mid q$, then $q^{30} \equiv 0 \pmod{5}$; otherwise Fermat's little theorem gives $q^{30} = (q^4)^7 q^2 \equiv q^2 \in \{1, 4\} \pmod{5}$. Thus $n \in \{1, 2, 3\} \pmod{5}$, so $5 \nmid n$.
\item Modulo $7$: $q^{30} = (q^6)^5 \in \{0, 1\} \pmod{7}$, so $n \in \{2, 3\} \pmod{7}$ and $7 \nmid n$.
\item Modulo $11$: $q^{30} = (q^{10})^3 \in \{0, 1\} \pmod{11}$, so $n \in \{2, 3\} \pmod{11}$ and $11 \nmid n$.
\item Modulo $13$: if $13 \mid q$, then $q^{30} \equiv 0 \pmod{13}$; otherwise $q^{12} \equiv 1 \pmod{13}$ implies $(q^6)^2 \equiv 1 \pmod{13}$, hence $q^{30} = (q^{12})^2 q^6 \equiv q^6 \in \{1, 12\} \pmod{13}$. Thus $n \in \{1, 2, 3\} \pmod{13}$, so $13 \nmid n$.
\end{itemize}
From $4 \nmid n$ we obtain $4, 8, 12 \nmid n$; from $5 \nmid n$ we obtain $5, 10 \nmid n$; together with $7, 9, 11, 13 \nmid n$, every positive divisor $d \mid n$ satisfying $d \le 13$ lies in $\{1, 2, 3, 6\}$. Therefore, if $S$ were a set of positive divisors of $n$ summing to $13$, every element of $S$ would be at most $13$, forcing $S \subseteq \{1, 2, 3, 6\}$ and
\[
13 = \sum_{d \in S} d \le 1 + 2 + 3 + 6 = 12,
\]
a contradiction. Since $13 \le n$, $n$ is not practical.

\medskip
Finally, if $C = 30 t$ with $t \ge 1$, then $q^C + 2 = (q^t)^{30} + 2$ with $q^t \ge 1$, which is never practical by the $C = 30$ case.
\end{proof}

\section{Exactness of the Modulus $30$ and Small-Exponent Contrast}

In the Stewart--Sierpi\'{n}ski criterion, an even integer $n$ with $4 \nmid n$ and $9 \nmid n$ has $\{2, 3\}$-part dividing $6$, and since $\sigma(6) + 1 = 13$, practicality requires a prime factor in $\{5, 7, 11, 13\}$. The next proposition shows that $30 \mid C$ is the exact characterization of this local obstruction.

\begin{proposition}\label{prop:sharp}
Let $C \ge 1$ be an integer. The local obstruction
\[
\forall q \in 2\mathbb{N},\qquad 4 \nmid q^C + 2,\; 9 \nmid q^C + 2,\; 5 \nmid q^C + 2,\; 7 \nmid q^C + 2,\; 11 \nmid q^C + 2,\; 13 \nmid q^C + 2
\]
holds if and only if $30 \mid C$.
\end{proposition}

\begin{proof}
Sufficiency follows from the proof of Theorem~\ref{thm:main} applied to $q^t$ (which is even whenever $q$ is even and $t \ge 1$). For necessity, suppose the obstruction holds for $C \ge 1$:
\begin{enumerate}
\item If $2 \nmid C$, then $C \equiv 1$ or $3 \pmod{4}$. Since $8^4 \equiv 2^4 \equiv 1 \pmod{5}$, taking $q = 8$ when $C \equiv 1 \pmod{4}$ gives $8^C + 2 \equiv 8 + 2 \equiv 0 \pmod{5}$, and taking $q = 2$ when $C \equiv 3 \pmod{4}$ gives $2^C + 2 \equiv 8 + 2 \equiv 0 \pmod{5}$, contradicting $5 \nmid q^C + 2$. Hence $2 \mid C$.
\item Since $2 \mid C$, we have $C \equiv 0, 2$, or $4 \pmod{6}$. As $4^6 \equiv 2^6 \equiv 1 \pmod{9}$, taking $q = 4$ for $C \equiv 2 \pmod{6}$ ($4^2 + 2 = 18$) or $q = 2$ for $C \equiv 4 \pmod{6}$ ($2^4 + 2 = 18$) gives $9 \mid q^C + 2$. Hence $3 \mid C$.
\item Since $2 \mid C$, we have $C \equiv 0, 2, 4, 6$, or $8 \pmod{10}$. As $q^{10} \equiv 1 \pmod{11}$ for $11 \nmid q$, the even witnesses $q = 8, 16, 2, 4$ satisfy $q^C \equiv 9 \pmod{11}$ (hence $11 \mid q^C + 2$) for $C \equiv 2, 4, 6, 8 \pmod{10}$, respectively. Hence $5 \mid C$.
\end{enumerate}
Combining $2 \mid C$, $3 \mid C$, and $5 \mid C$ yields $30 \mid C$.
\end{proof}

By contrast, for every small exponent $1 \le C \le 6$ there exist explicit practical numbers $q$ for which both $q^C$ and $q^C + 2$ are practical:
\[
(q, C) \in \{(2, 1),\, (2, 2),\, (4, 3),\, (2, 4),\, (4, 5),\, (2, 6)\},
\]
yielding $q^C + 2 \in \{4, 6, 66, 18, 1026, 66\}$, all of which are practical. Thus the conjecture remains open for positive exponents $C$ not divisible by $30$ outside the previously proved cases ($C \in \{1, 2, 4\}$ and primes $7 \le C < 100$).

\section{Formal Verification in Lean~4}

All results above are formally machine-checked in Lean~4 (\texttt{v4.33.1} + Mathlib) in the module \texttt{LiTwinPractical.lean}. Every theorem depends only on the three standard axioms (\texttt{propext}, \texttt{Classical.choice}, and \texttt{Quot.sound}), with zero \texttt{sorry}, \texttt{admit}, or \texttt{native\_decide}:
\begin{itemize}
\item Theorem~\ref{thm:main}: \texttt{not\_practical\_pow30\_add\_two}, \texttt{not\_practical\_pow\_mul30\_add\_two},\\
\texttt{li\_conjecture\_1\_1\_empty\_of\_dvd\_thirty}, and \texttt{li\_conjecture\_1\_1\_false};
\item Proposition~\ref{prop:sharp}: \texttt{local\_obstruction\_iff\_dvd\_thirty};
\item Small-exponent positive witnesses ($1 \le C \le 6$): \texttt{small\_exponents\_witnesses}.
\end{itemize}

\begin{thebibliography}{4}
\bibitem{Li2026}
R. Li,
\emph{On a conjecture involving twin practical numbers},
\emph{Notes on Number Theory and Discrete Mathematics} \textbf{32}(3) (2026), 544--555,
doi:10.7546/nntdm.2026.32.3.544-555.

\bibitem{Sierpinski1955}
W. Sierpi\'{n}ski,
\emph{Sur une propri\'{e}t\'{e} des nombres naturels},
\emph{Ann. Mat. Pura Appl.} \textbf{39} (1955), 69--74.

\bibitem{Srinivasan1948}
A. K. Srinivasan,
\emph{Practical numbers},
\emph{Current Science} \textbf{17} (1948), 179--180.

\bibitem{Stewart1954}
B. M. Stewart,
\emph{Sums of distinct divisors},
\emph{Amer. J. Math.} \textbf{76} (1954), 779--785.

\bibitem{WangSun2018}
L. Wang and Z.-W. Sun,
\emph{On practical numbers of the form $a n^2 + b n + c$},
arXiv:1809.01532 (2018).
\end{thebibliography}

\end{document}
